%\documentclass[11pt,a4paper,twoside]{article}

\usepackage[utf8]{inputenc}

\usepackage[T1,T2A]{fontenc}

\usepackage[english.ancient,russian.ancient]{babel}

\usepackage{graphicx}

\usepackage{array}

\usepackage[doi]{samgtu-bib}

\usepackage{samgtu}

\usepackage{samgtu-name}

\usepackage{mathtools}

\mathtoolsset{showonlyrefs}

\usepackage{desclist}

\usepackage{euscript}

\usepackage{mathrsfs} 

\usepackage{multicol}

\usepackage{mathrsfs}

\usepackage{cite}

\usepackage{cmap}

\usepackage{qrcode}

\allowdisplaybreaks[4]

\usepackage[nodayofweek]{datetime}

\usepackage[thinlines]{easybmat}



\renewcommand{\JournalYear}{2018}

\renewcommand{\JournalVolume}{22}

\renewcommand{\JournalNumber}{4}



\newdate{JournalDateSign}{29}{12}{2018}% Дата подписания Вестника в печать

\renewcommand{\JournalDateSign}{\displaydate{JournalDateSign}}

\newdate{JournalDatePrint}{16}{01}{2019}% Дата выхода Вестника из печати

\renewcommand{\JournalDatePrint}{\displaydate{JournalDatePrint}}

%\renewcommand{\JournalUPL}{16.56} % Усл. печ. л.

%\renewcommand{\JournalUIL}{16.53} % Уч.-изд. л.

\renewcommand{\JournalUPL}{15.24} % Усл. печ. л.

\renewcommand{\JournalUIL}{15.21} % Уч.-изд. л.

\renewcommand{\JournalRegNumber}{220/18} % Регистрационный номер

\renewcommand{\JournalOrderNumber}{2} % Номер заказа



%\renewcommand{\JournalRMonth}{Март} % Месяц выхода Вестника

%\renewcommand{\JournalEMonth}{March} % Месяц выхода Вестника

%\renewcommand{\JournalRMonth}{Июнь} % Месяц выхода Вестника

%\renewcommand{\JournalEMonth}{June} % Месяц выхода Вестника

%\renewcommand{\JournalRMonth}{Сентябрь} % Месяц выхода Вестника

%\renewcommand{\JournalEMonth}{September} % Месяц выхода Вестника

\renewcommand{\JournalRMonth}{Декабрь} % Месяц выхода Вестника

\renewcommand{\JournalEMonth}{December} % Месяц выхода Вестника



\newcommand{\totop}[1]{\raisebox{6pt}[1pt][2pt]{#1}}

\newcommand{\tottop}[1]{\raisebox{12pt}[1pt][2pt]{#1}} 

\newcommand{\totttop}[1]{\raisebox{18pt}[1pt][2pt]{#1}} 

\newcommand{\tottttop}[1]{\raisebox{24pt}[1pt][2pt]{#1}} 

\newcommand{\totttttop}[1]{\raisebox{30pt}[1pt][2pt]{#1}} 

\newcommand{\tobot}[1]{\raisebox{-6pt}[1pt][2pt]{#1}} 

\newcommand{\tobbot}[1]{\raisebox{-12pt}[1pt][2pt]{#1}} 

\newcommand{\tobbbot}[1]{\raisebox{-18pt}[1pt][2pt]{#1}}



\graphicspath{{./00PIC/}}

%

%\usepackage{samgtu-name}

%

\vsgtu{1642} %registration number

\FirstPage{574}

\LastPage{585}

%

\ShortCommunication

%

%\pdfcrop

%\draft

%

%\begin{document}

%

%

%$\;$

%

%\hfill \qrcode[hyperlink,height=.8in]{http://doi.org/10.14498/vsgtu1642}

%\vspace{-.8in}





\renewcommand{\baselinestretch}{0.88}



\udc{539.422.5:539.4.01}

\msc{74A40, 82D80}



\rutitle{Исследование прогиба  композита <<УНТ-графен>> с~использованием молекулярно-динамического моделирования}

\rutitleshort{Исследование прогиба  композита <<УНТ-графен>>\dots}



\entitle{Investigation of deflection of the CNT/G composite by~molecular dynamics simulation}





\ruauthor{А.~С.~Колесникова, И.~В.~Кириллова, \\ 

Г.~А.~Барегамян, Л.~Ю.~Коссович}{К\,о\,л\,е\,с\,н\,и\,к\,о\,в\,а~А.~С., К\,и\,р\,и\,л\,л\,о\,в\,а~И.~В., Б\,а\,р\,е\,г\,а\,м\,я\,н~Г.~А., К\,о\,с\,с\,о\,в\,и\,ч~Л.~Ю.}

\enauthor{A.~S.~Kolesnikova, I.~V.~Kirillova, \\ 

G.~A.~Baregamyan, L.~Yu.~Kossovich}{K\,o\,l\,e\,s\,n\,i\,k\,o\,v\,a~A.~S., K\,i\,r\,i\,l\,l\,o\,v\,a~I.~V., B\,a\,r\,e\,g\,a\,m\,y\,a\,n~G.~A., K\,o\,s\,s\,o\,v\,i\,c\,h~L.~Yu.}



\ruaffil{\saratovuni.}



\enaffil{\engsaratovuni.}



\ruauthorsdetails{4}{% Количество авторов

\fullauthor{\rupid{79039}{Анна Сергеевна Колесникова}}

\orcid{0000-0003-2033-1192} % ORCID ID автора 

\regalia{кандидат физико-математических наук} 

\position{доцент} 

\dept{каф. математической теории упругости и~биомеханики} 

\email{Kolesnikova.88@mail.ru}

\fullauthor{\rupid{65717}{Ирина Васильевна Кириллова}}

\CAuthor % Автор-корреспондент

\orcid{0000-0001-8053-3680} % ORCID ID автора 

\regalia{кандидат физико-математических наук} 

\position{директор} 

\dept{научно образовательный институт наноструктур и биосистем} 

\email{ivkirillova@yandex.ru}

%\fullauthor{Гайк Барегамян}

%%\orcid{0000-0000-0000-xxxx} % ORCID ID автора 

%\position{аспирант} 

%\dept{каф. математической теории упругости и биомеханики} 

%\email[br]{baregaik@gmail.com}

%\fullauthor{Леонид Юрьевич Коссович}

%\orcid{0000-0002-4775-7348} % ORCID ID автора 

%\regalia{доктор физико-математических наук} 

%\position{заведующий кафедрой} 

%\dept{каф. математической теории упругости и биомеханики} 

%\email{president@sgu.ru}

}



\enauthorsdetails{4}{

\fullauthor{\enpid{79039}{Anna S. Kolesnikova}}

\orcid{0000-0003-2033-1192} % ORCID ID

\regalia{Cand. Phys. \& Math. Sci.} 

\position{Associate Professor} 

\dept{Dept. of Mathematical Theory of Elasticity and Biomechanics}

\email{Kolesnikova.88@mail.ru}

\fullauthor{\enpid{65717}{Irina V. Kirillova}}

\CAuthor % Сorresponding Author

\orcid{0000-0001-8053-3680} % ORCID ID

\regalia{Cand. Phys. \& Math. Sci.} 

\position{Director} 

\dept{Scientific and Educational Institute of Nanostructures and Biosystems}

\email{ivkirillova@yandex.ru}

%\fullauthor{G.~Baregamyan}

%%\orcid{0000-0000-0000-xxxx} % ORCID ID

%\position{Postgraduate Student} 

%\dept{Dept. of mathematical theory of Elasticity and Biomechanics}

%\email{baregaik@gmail.com}

%\fullauthor{Leonid Yu. Kossovich}

%\orcid{0000-0002-4775-7348} % ORCID ID avtora 

%\regalia{Dr. Phys. \& Math. Sci., Professor} 

%\position{Head of Department} 

%\dept{Dept. of mathematical theory of Elasticity and Biomechanics}

%\email{president@sgu.ru}

}



\rukeywords{прогиб, углеродный композит, молекулярно-ди\-на\-ми\-чес\-кое моделирование, индентирование}

\enkeywords{deflection, carbon composite, molecular dynamics simulation, indentation}



\ruabstract{В последнее время графен и наноматериалы на его основе используются в области биомедицины в качестве материала для биосенсорики. Основными компонентами в биосенсорах являются датчики, которые должны быть гибкими, масштабируемыми, чувствительными и надежными. Деформация материала изменяет его электрическое сопротивление, поэтому исследование механических свойств композитов, состоящих из нанотрубок и графена, является важной задачей. 



В настоящее время осуществляется активная разработка методик синтеза композитов, состоящих из графена и параллельно ему ориентированных нанотрубок. В связи с этим в данной работе проведено теоретическое исследование зависимости изгибающей силы от поперечного смещения атомов в центре композитного материала, состоящего из графена и параллельно ему ориентированных $(8,0)$ zigzag-нанотрубок. Выбор нанотрубки с~хиральностью  $(8,0)$ для исследования в данной работе обусловлен минимальным диаметром нанотрубок, входящих в состав композита данного типа. 

Устойчивость композита оценивалась путем расчета по величине энтальпии и связывалась с отрицательными ее значениями. Установлено, что энтальпия не изменяется в зависимости от расстояния между осями, вдоль которых ориентированы нанотрубки, входящие в~состав композитов.



Композитный материал удерживается с~двух краев опорами при отсутствии подложки. Поиск равновесного состояния структуры определялся методом молекулярной механики с использованием энергетического потенциала Бреннера в~рамках метода молекулярной динамики. 



Математическое моделирование действия иглы атомно"=силового микроскопа осуществлялось при помощи однослойной  углеродной нанотрубки типа armchair. Взаимодействие между armchair-нанотрубкой и~композитом осуществляется за счет сил Ван--дер--Ваальса. 

Нанотрубки типа armchair удобны для их использования в качестве иглы атомно-силового микроскопа, потому что их край не имеет острых углов (в отличие от трубок типа zigzag), т.~е. при контакте будут наноситься минимальные повреждения системе.



Приводятся результаты расчетов для зависимости энтальпии от диаметра нанотрубок, стрелы прогиба в центре пластины из композитного материала <<УНТ-графен>> от приложенной силы при изгибе. Показано, что исследуемые в данной работе нанокомпозиции энергетически устойчивы.}



\enabstract{Graphene and nanomaterials based on graphene have been using in the field of biomedicine as a material for biosensorics. The main components in biosensors are sensors, which must be flexible, scalable, sensitive and reliable. The deformation of the material changes its electrical resistance, therefore the study of the mechanical properties of composites, consisting of nanotubes and graphene, is the urgent task. 



Currently, active development of methods for the synthesis of composites consisting of graphene and parallel to it oriented nanotubes have been carrying. However, papers on the investigation of the optical and electronic properties of this composition was carried out not enough, and papers on the investigation of the mechanical properties of composites have not been found. The aim of this work is a theoretical investigation of the depending the bending force on the transverse displacement of atom in center of the composite material consisting of graphene and parallel to it $(8,0)$ zigzag nanotubes. The choice of a nanotube $(8,0)$ for research in this work is due to the minimum diameter of the nanotubes that make up the composite of this type. The stability of the composite was estimated by calculating the value of enthalpy and is characterized by a negative value of enthalpy. It was established that enthalpies do not change depending on the distance between the axes, along which the nanotubes belonging to the composites are oriented. 



Composite material was retained on both edges by support in the absence of a substrate. The search for the equilibrium state of the structure was determined by the molecular mechanics method using the Brenner energy potential within the framework of the molecular dynamics method. 



Mathematical modeling of the action of the needle of the atomic force microscope was carried out using the single-layer armchair carbon nanotube. The interaction between the armchair nanotube and the composite is carried out by means of the van der Waals forces.}



\newdate{kolesnikovaa}{18}{08}{2018}

\date{\displaydate{kolesnikovaa}}

\newdate{kolesnikovab}{28}{08}{2018}

\revisiondate{\displaydate{kolesnikovab}}

\newdate{kolesnikovac}{03}{09}{2018}

\accepteddate{\displaydate{kolesnikovac}}



\newdate{kolesnikovae}{08}{10}{2018}

\renewcommand{\JournalDataOnline}{\displaydate{kolesnikovae}}



%\ForPeerReview

\makerutitle



\footnotetext{\fullauthor{\rupid{73316}{Гайк Артурович Барегамян}}

%\orcid{0000-0000-0000-xxxx} % ORCID ID автора 

\position{аспирант} 

\dept{каф. математической теории упругости и биомеханики} 

\email{baregaik@gmail.com}

\fullauthor{\rupid{14470}{Леонид Юрьевич Коссович}}

\orcid{0000-0002-4775-7348} % ORCID ID автора 

\regalia{доктор физико-математических наук} 

\position{заведующий кафедрой} 

\dept{каф. математической теории упругости и биомеханики} 

\email{president@sgu.ru}}





\Section[n]{Введение}

Перспектива развития наноустройств связана с нахождением новых уникальных свойств углеродных материалов, 

что стимулирует синтез и исследование одного из перспективных новых материалов, особенности и возможности которых не были раскрыты в полной мере. Этим материалом является углеродный композит <<УНТ-графен>>, образованный углеродными нанотрубками (УНТ) и~графеном.



Интерес именно к этому композиту вызван тем, что графен является углеродным материалом с высокой электрической проводимостью, теплопроводностью, оптической прозрачностью и исключительными механическими свойствами \cite{kol:bib1}. В настоящее время отсутствует технология его массового изготовления из-за сложности получения графена протяженных размеров, что связано с наличием дефектов в синтезируемых образцах графена. Поэтому для решения этой проблемы предложено синтезировать новую гибридную структуру, состоящую из ковалентно соединенных между собой однослойных углеродных нанотрубок (УНТ) и графеновых листов \cite{kol:bib2}.



В настоящее время синтезируемые композиты <<УНТ-графен>> можно разделить на две морфологические группы \cite{kol:bib3}: 

\begin{itemize}

\item[1)] 3D-композит "--- <<колонный графен>> "--- нанотрубки ориентированы вертикально относительно графеновых слоев и расположены между ними; 

\item[2)] 2D гибридные пленки "--- трубки расположены между соседними слоями, параллельными им (ковалентное или ван-дер-ваальсовое соединение между УНТ и~графеном). 

\end{itemize}

Во второй группе можно выделить структуры с упорядоченным и хаотическим параллельным расположением УНТ друг относительно друга.



Среди двух классов композитов <<УНТ-графен>> особый интерес вызывает класс композитов 2D гибридные пленки. 

Этот интерес обусловлен тем, что научных работ по исследованию свойств композита 2D гибридные пленки мало \cite{kol:bib4,kol:bib5,kol:bib6,kol:bib7,kol:bib8}. 

В~работе \cite{kol:bib5}  осуществлялось исследование зависимости электрического сопротивления композита <<УНТ-графен>> от величины деформации изгиба.  

Авторы работ \cite{kol:bib4,kol:bib5}  показали, что электрическое сопротивление композита <<УНТ-графен>> ниже по сравнению с~электрическим сопротивлением составных частей этого композита, т.~е. углеродных нанотрубок и~графена. 

Композиты типа <<УНТ-графен>> уже используются в ультрапрозрачных гибких устройствах памяти \cite{kol:bib9}, в~полевых транзисторах \cite{kol:bib4}, а также в качестве суперконденсаторов~\cite{kol:bib10} и~материала для электрохимического хранения энергии \cite{kol:bib11,kol:bib12}. 

Таким образом, можно отметить, что, несмотря на большое разнообразие потенциальных применений гибридного <<УНТ-графен>> "--- композита нового поколения, направление его исследования, связанное с механическими свойствами, не изучено.



Для расширения областей применения такого композита важной задачей является подробное исследование его свойств. Однако основной задачей улучшения и модернизации наноустройств путем использования в~них композитов <<УНТ-графен>> в~качестве элементной базы необходимо исследовать механические свойства этих структур. 

В~настоящее время исследование данных характеристик не проводилось для композита <<УНТ-графен>> в~виде 2D гибридных пленок.



Целью данной работы является теоретическое исследование динамической зависимости изгибающей силы от стрелы прогиба композитного материала <<УНТ-графен>>, имеющего вид 2D гибридных пленок. 

Композит удерживался с~двух краев опорами при отсутствии подложки. 

Поиск равновесного состояния структуры определялся методом молекулярной механики с~использованием энергетического потенциала Бреннера \cite{kol:bib13} в~рамках метода молекулярной динамики \cite{kol:bib14,kol:bib15,kol:bib16}.



\smallskip



\Section{Объект исследования}

Объектом исследования является углеродный композит, образованный двумя графеновыми лентами типа zigzag и тремя нанотрубками типа zigzag, расположенными между ними и соединенными с~ними ковалентными химическими связями (рис.~\ref{fig1}). 

В~целом композит представляет собой 2D-пленку, протяженную в направлениях $X$ и $Y$. 

Нанотрубки в композите ориентированы параллельно друг другу и параллельно графеновой ленте. 

Выбор типа хиральности нанотрубок в~композите и~способ расположения нанотрубок относительно графеновой ленты осуществляется на основе структурной модели композита <<УНТ-графен>>, представленной в~экспериментальной работе \cite{kol:bib7}.



\begin{figure}[p]

\center  \footnotesize 

\includegraphics[scale=.6]{_fig1}

\caption{Композит <<УНТ-графен>> с~параллельным расположением УНТ относительно графеновых лент:  

$d$ "--- расстояние между осями смежных нанотрубок; $L$ "--- длина графеновой наноленты типа zigzag; 

$W$ "--- ширина  графеновой наноленты типа armchair; $h$ "--- расстояние между двумя графеновыми нанолентами}

\label{fig1}

\smallskip \footnotesize

[Figure~\ref{fig1}. The structure of the CNT/G composite with parallel arrangement of CNT relative to graphene ribbons; $d$ is distance between the axes of adjacent nanotubes; $L$ is the length of the zigzag graphene nanoribbons; $W$ is the width of the armchair of graphene nanoribbons; $h$ is the distance between two graphene nanoribbons]



\bigskip



\center 

%\includegraphics[width = 15cm, length= 27cm]{fig2}

\includegraphics[width = 0.8\textwidth]{_fig2}

\includegraphics[width = 0.6\textwidth]{_fig3}

\caption{Фрагменты композитов <<УНТ-графен>>, содержащие углеродные $(9,0)$ zigzag"=нанорубки (сверху слева),  $(10,0)$ zigzag-нанорубки (сверху справа) и  $(6,0)$ zigzag-нанорубки (снизу)}

\label{fig2}

\smallskip \footnotesize

[Figure~\ref{fig2}. Fragments of the structure of the CNT/G composite containing  $(9,0)$ zigzag nanotube (top left), $(10,0)$ zigzag nanotube (top right), and $(6,0)$ zigzag nanotube (bottom)]

\end{figure}



В данной работе осуществляется исследование композитов, в состав которых входят zigzag"=нанотрубки. 

Отметим, что выбирались нанотрубки с хиральностью $(n,0)$, где $n$ "--- четное число, чтобы количество химических связей нанотрубок с~лентами (нижней и верхней) было одинаковым "--- равным двум. 

Иначе, как показано на рис.~\ref{fig2}, количество связей может быть разным. 

На основе результатов экспериментальных данных известно, что стабильными zigzag"=нанотрубками наименьшего диаметра являются $(6,0)$ zigzag"=нанотрубки \cite{kol:bib17}. 

Однако при взаимодействии графеновой ленты с $(6,0)$ zigzag"=нанотрубкой наблюдается уменьшение количества химических связей между УНТ и графеном по сравнению с количеством химических связей, образующихся при взаимодействии нанотрубок с хиральностью $(8,0)$, $(10,0)$, $(12,0)$, $(14,0)$, $(16,0)$ с графеновыми нанолентами (рис.~\ref{fig2}). 

Это объясняется увеличением угла пирамидализации атомного каркаса $(6,0)$ zigzag"=нанотрубки по сравнению с углом пирамидализации zigzag"=нанотрубок с хиральностью $(8,0)$, $(10,0)$, $(12,0)$, $(14,0)$, $(16,0)$. 

В~связи с этим выбор нанотрубки с хиральностью  $(8,0)$ для исследования в~данной работе обусловлен минимальным диаметром нанотрубок, входящих в~состав композита данного типа. 

При дальнейшем увеличении диаметра УНТ будет наблюдаться формирование композитов с~данной конфигурацией, т.~е. с~таким количеством химических связей между графеном и~УНТ.



%\begin{figure}[t!]

%\center 

%%\includegraphics[width = 15cm, length= 27cm]{fig2}

%\includegraphics[width = 0.8\textwidth]{fig2}

%\includegraphics[width = 0.6\textwidth]{fig3}

%\caption{Фрагменты композитов <<УНТ-графен>>, содержащих углеродные $(9,0)$ zigzag-нанорубки (слева) и $(10,0)$ zigzag-нанорубки (справа) }

%\label{fig2}

%\smallskip \footnotesize

%[Figure~\ref{fig2}. Fragments of the structure of the CNT/G composite containing  $(9,0)$ zigzag nanotube (left) and  $(10,0)$ zigzag nanotube (right)]

%\end{figure}

%

%\begin{figure}[h]

%\center 

%%\includegraphics[width = 8 cm, length= 27cm]{fig3}

%\includegraphics[width = 0.6\textwidth]{fig3}

%\caption{Фрагмент композита <<УНТ-графен>>, содержащий zigzag $(6,0)$ углеродные нанорубки.}

%\label{fig3}

%\smallskip \footnotesize

%[Figure~\ref{fig3}. Fragment of the composite <<CNT graphene>> containing zigzag (6,0) carbon nanotubes.]

%\end{figure}



Длина нанотрубок и ширина края графеновой armchair"=ленты   являются фиксированными и~составляют 2.84~нм. 

Длина края графеновой zigzag"=ленты  составляет 5.88~нм. 

Исследование энергетической устойчивости композитных конструкций характеризуется отрицательным значением энтальпии. Энтальпия рассчитывается следующим образом \cite{kol:bib18}:

\begin{equation}

\label{kol:1}

 E=(E_{tot}-E_{1-5}-E_{tube})/N_{atom},%\; \rm  [eV/atom],

\end{equation}

где $E_{tot}$ "--- полная энергия равновесного состояния построенной модели композита <<УНТ-графен>>, рассчитанная с~использованием модифицированного эмпирического метода, основанного на энергетическом потенциале Бреннера \cite{kol:bib13}; 

$E_{1-5}$ "--- полная энергия составных частей композита <<УНТ-графен>>, находящихся в равновесном состоянии (два графеновых листа и три углеродные нанотрубки); 

$E_{tube}$ "--- суммарная энергия нанотрубок, из которых состоит композит; $N_{atom}$ "--- количество атомов в~композите <<УНТ-графен>>. 

Полная энергия каждой составной части композита <<УНТ-графен>> также рассчитывается с использованием модифицированного эмпирического метода, основанного на энергетическом потенциале Бреннера. 



\smallskip



\Section{Результаты исследования}

Проведено исследование зависимости энтальпии композита при изменении диаметра УНТ, которые входят в состав композита, и при увеличении расстояния между осями, вдоль которых ориентированы нанотрубки в композите. 

Исследования проводились для трех расстояний между осями смежных УНТ, вдоль которых эти УНТ ориентированы в~композите: 1.46~нм,   1.96~нм  и~2.21~нм. 

Диаметры нанотрубок менялись от 0.24~нм до 0.65~нм. 

В~процессе исследований установлено, что энтальпия не изменяется в~зависимости от расстояния между осями, вдоль которых ориентированы нанотрубки, входящие в~состав композитов. 

Например, если диаметр zigzag"=нанотрубки составляет 0.39~нм, то значение энтальпии составляет минус $2.38$~эВ/атом для трех разных расстояний между осями, вдоль которых ориентированы нанотрубки. 

Повысить энтальпию нанокомпозита <<УНТ-графен>> можно за счет уменьшения диаметра нанотрубок, входящих в~состав этого композита. 

На рис. \ref{fig4} приведен график зависимости энтальпии от диаметра нанотрубок, которые входят в~состав композита <<УНТ-графен>>, при расстоянии между осями смежных УНТ, вдоль которых ориентированы нанотрубки, равном 1.96~нм.



\begin{figure}[h!]

\center 

\begin{tabular}{cp{.48\textwidth}}



\includegraphics[ scale=0.3]{_fig4} & \vspace{-6cm}

\caption{Зависимость энтальпии от диаметра нанотрубок, которые входят в состав композита <<УНТ-графен>>, при расстоянии между осями смежных УНТ, вдоль которых ориентированы нанотрубки, равном 1.96~нм

\label{fig4}}



\smallskip \footnotesize

[Figure~\ref{fig4}. The dependence of the enthalpy on the diameter of CNTs which are part of the CNT/G composite, when the distance 1.96~\textmu m between the axes of adjacent CNTs along \centerline{ which the nanotubes are oriented]}



\\



\end{tabular}



\vspace{-5mm}



\end{figure}



Процесс изгиба осуществлялся на композите с диаметром нанотрубок \linebreak 0.61~нм,  длиной нанотрубок 2.87~нм и~расстоянием между нанотрубками  1.96~нм. 

Значение энтальпии данной структуры составляет минус $2.1$ эВ/атом, что свидетельствует о ее стабильности. Исследование процесса изгиба проводилось путем воздействия на композит иглы атомно"=силового микроскопа, как это осуществляется в~методе наноиндентирования. 

Однако данное исследование отличается от метода наноиндентирования тем, что в~методе наноиндентирования имеется подложка, на которой располагается образец под действием силы со стороны иглы наноиндентора, а в~данном исследовании подложка отсутствует, т.~к. образец удерживается двумя опорами, располагающимися по бокам от образца (рис.~\ref{fig5}).



\begin{figure}[b!]



\vspace{-3mm}



\center 

\begin{tabular}{cp{.48\textwidth}}



\includegraphics[width = 0.45\textwidth]{_fig5} & \vspace{-2.5cm}

\caption{Композит, образованный двумя графеновыми лентами и тремя нанотрубками, удерживающийся с помощью двух опор

\label{fig5}}



\smallskip \footnotesize

[Figure~\ref{fig5}. The CNT/G composite formed by two graphene ribbons and three nanotubes \centerline{held by 2 supports]}



\\



\end{tabular}









\end{figure}



Математическое моделирование действия иглы атомно-силового микроскопа осуществлялось при помощи однослойной  углеродной armchair"=нанотрубки. 

Взаимодействие между armchair"=нанотрубкой и~композитом осуществляется за счет сил Ван--дер--Ваальса. 

Нанотрубки типа armchair удобны для использования в~качестве иглы атомно"=силового микроскопа, так как их край не имеет острых углов (в отличие от трубок типа zigzag) и~они будут наносить минимальное повреждение системе при контакте. 

Аналогичная нанотрубка была выбрана для исследования процесса индентирования липопротеина \cite{kol:bib19}. 

Размеры нанотрубки были следующими: длина УНТ составила 2.97~нм, диаметр трубки "--- 1.24~нм. 

При моделировании атомы углеродной нанотрубки жестко закреплялись, чтобы она не деформировалась в~процессе исследования. 

Игла в~виде тонкой одностенной углеродной нанотрубки опускалась сверху вниз, перпендикулярно к графену, приблизительно по центру композита. 

Чтобы композиту не дать возможность уйти от иглы, т.~е. ограничить число степеней свободы, атомы композита на торце armchair"=ленты жестко закреплялись, моделируя тем самым опоры с~двух сторон композита. 

Расстояние между структурой и нижним сечением трубки "--- 0.33~нм (равновесное физическое расстояние).



На каждом шаге, через 500~фсек, трубка приближалась к композиту на $\xi=0.02$~нм. 

На каждом шаге фиксировались значения прогиба $\delta$ в~центре пластины. 

Величина $\delta$ находилась как разность координат атомов (по оси~$z$) между начальными и~конечными значениями соответственно (рис.~\ref{fig6}). 

За нулевой уровень прогиба принято положение композита в~первый момент контакта с~углеродной нанотрубкой (см. рис.~\ref{fig4}).



\begin{figure}[h!]

\center 

%\includegraphics[width = 9 cm, length= 27cm]{fig6}

\includegraphics[width = 0.6\textwidth]{_fig6}

\caption{Изогнутый композит при стреле прогиба  $\delta=0.66$~нм

\label{fig6}}

\smallskip \footnotesize

[Figure~\ref{fig6}. Curved composite with deflection of $\delta=0.66$ \textmu m]

\end{figure}





\begin{figure}[b!]



\vspace{-3mm}

\center 

\begin{tabular}{cp{.48\textwidth}}



\includegraphics[width = 0.4\textwidth]{_fig7} & \vspace{-2cm}

\caption{График зависимости приложенной силы $F$ от стрелы прогиба $\delta$ в центре пластины

\label{fig7}}



\smallskip \footnotesize

[Figure~\ref{fig7}. The plot of the applied force $F$ on the deflection $\delta$  in the center of the plate]

\\



\end{tabular}



\vspace{-3mm}



\end{figure}



Расчет силы $F$, необходимой для прогиба композита $\delta$, осуществлялся с~помощью подхода, предложенного в работе \cite{kol:bib20}: 

\begin{equation}

  F=W/\zeta,

\end{equation}

где $W$ "--- энергия взаимодействия, в качестве которой берется потенциал Леннарда"--~Джонса, 

$\zeta$ "--- значение прогиба на каждом шаге.



Увеличение энергии взаимодействия между нанотрубкой и~композитом вызвано ван-дер-ваальсовым отталкиванием между нанотрубкой и~композитом. 

График зависимости приложенной силы $F$ от стрелы прогиба $\delta$ в~центре пластины изображён на рис.~\ref{fig7}. Из графика видно, что стрела прогиба композита увеличивается с~увеличением силы воздействия. 

Установлено, что силы 6~мкН недостаточно для того, чтобы наблюдался процесс разрушения композита. 

Об этом свидетельствует изменение длин химических связей в~центральной области композита, на которую оказывается наибольшее воздействие со стороны нанотрубки "--- иглы наноиндентора. 

Эти длины связей составляют в~среднем 0.14~нм, а~разрыв химической связи наблюдается при~0.19~нм.











\Section[n]{Выводы}

Впервые выполнено исследование динамической зависимости изгибающей силы от стрелы прогиба композитного материала <<УНТ-графен>>. 

Исследование процесса изгиба композита проводилось путем моделирования воздействия на композит иглы атомно"=силового микроскопа, представляющей собой однослойную нанотрубку суб- и нанометрового диаметров. 

Композит имел следующую структуру: углеродный материал, образованный двумя графеновыми zigzag"=лентами и тремя zigzag"=нанотрубками, расположенными между ними и соединенными с~ними ковалентными химическими связями. 

В~рамках исследования установлено, что минимальный диаметр нанотрубок, входящих в~состав композита <<УНТ-графен>>, составляет 0.615~нм. При уменьшении диаметра УНТ в~композите наблюдается уменьшение количества химических связей между нанотрубками и~графеновой лентой, такая конфигурация отмечается во вcех композитах при увеличении диаметра УНТ типа zigzag.



В результате численного эксперимента установлено, что исследуемые в~данной работе нанокомпозиты <<УНТ-графен>> являются энергетически устойчивыми. 

Для исследования динамической зависимости изгибающей силы от стрелы прогиба композитного материала <<УНТ-графен>> содержал нанотрубки диаметром 0.61~нм  и длиной 2.87~нм c~расстоянием между ними  1.96~нм. Ширина края графеновой armchair"=ленты является фиксированной и~составляет 2.84~нм. 

Длина края графеновой zigzag"=ленты составляет 5.88~нм. 

В~рамках данного исследования установлено, что при приложении к~композиту силы 6~мкН разрушения композита наблюдаться не будет.



Принимая во внимание существующие технологии успешного синтеза композитных материалов <<УНТ-графен>> \cite{kol:bib4,kol:bib5,kol:bib6}, можно предположить, что представленные в~данной работе результаты дадут возможность расширить область применения композитов <<УНТ-графен>>.



\AdditionalContent{Конкурирующие интересы}{Мы заявляем, что у нас нет конфликта интересов в авторстве и публикации этой статьи.} 



\AdditionalContent{Авторский вклад и ответственность}{Все  авторы  принимали  участие  в  разработке  концепции  \mbox{статьи}  и  в  написании  рукописи. Мы несем  полную  ответственность  за  предоставление  окончательной рукописи в печать. Окончательная  версия рукописи была нами одобрена.} 



\AdditionalContent{Финансирование}{Работа выполнена при финансовой поддержке Президентской стипендии 2016--2018 (проект № СП-2502.2016.1).}





\begin{thebibliography}{10}





\Bibitem{kol:bib1}

\by Wei~J., Wei~C., Su~L., Fu~J., Lv~J. 

\yr 2015

\paper Synergistic Reinforcement of Phenol-Formaldehyde Resin Composites by Poly(Hexanedithiol)/Graphene Oxide

\jour  J.~Mater. Sci. Chem. Eng.

\vol 3

\issue 8

\pages 56--70

\crossref{10.4236/msce.2015.38009}



\Bibitem{kol:bib2}

\by Potts~J.~R., Dreyer~D.~R., Bielawski~C.~W., Ruoff~R.~S.

\paper Graphene-based polymer nanocomposites

\jour Polymer

\yr 2011

\vol 52

\issue 1

\pages 5--25

\crossref{10.1016/j.polymer.2010.11.042}





\Bibitem{kol:bib3}

\by Jung~N., Kwon~S., Lee~D., Yoon~D.~M. et. al.

\paper Synthesis of Chemically Bonded Graphene/ Carbon Nanotube Composites and their Application in Large Volumetric Capacitance Supercapacitors

\jour Adv. Mater.

\yr 2013

\vol 25

\issue 47

\pages 6854--6858

\crossref{10.1002/adma.201302788}



\Bibitem{kol:bib4}

\by Kim~S.~H., Song~W., Jung~M.~W., Kang~M.~A. et. al.

\paper Carbon nanotube and graphene hybrid thin film for transparent electrodes and field effect transistors

\jour Adv. Mater.

\yr 2014

\vol 26

\issue 25

\pages 4247--4252

\crossref{10.1002/adma.201400463}





\Bibitem{kol:bib5}

\by Kholmanov~I.~N., Magnuson~C.~W. , Piner~R., Kim~J.~Y. et. al.

\paper Optical, electrical, and electromechanical properties of hybrid graphene/carbon nanotube films

\jour Adv. Mater.

\yr 2015

\vol 27

\issue 19

\pages 3053--3059

\crossref{10.1002/adma.201500785}





\Bibitem{kol:bib6}

\by Dong~X., Li~B.,  Wei~A., Cao~X. et. al.

\paper One-step growth of graphene–carbon nanotube hybrid materials by chemical vapor deposition

\jour Carbon

\yr 2011

\vol 49

\issue 9

\pages 2944--2949

\crossref{10.1016/j.carbon.2011.03.009}



\Bibitem{kol:bib7}

\by Trist\'an-L\'opez~F., Morelos-G\'omez~A., Vega-Diaz~S.~M., Garcia-Betancourt~M.~L. et. al.

\paper Large area films of alternating graphene-carbon nanotube layers processed in water

\jour ACS Nano

\yr 2013

\vol 7

\issue 12

\pages 10788--10798

\crossref{10.1021/nn404022m}





\Bibitem{kol:bib8}

\by Mitrofanov~V.~V., Slepchenkov~M.~M., Zhang~G., Glukhova~O.~E.

\paper Hybrid Carbon Nanotube–Graphene Monolayer Films: Regularities of Structure, Electronic and Optical Properties

\jour Carbon

\yr 2017

\vol 115

\pages 803--810

\crossref{10.1016/j.carbon.2017.01.040}





\Bibitem{kol:bib9}

\by Yu~W.~J., Chae~S.~H., Lee~S.~Y., Duong~D.~L., Lee~Y.~H.

\paper Ultra-transparent, flexible single-walled carbon nanotube non-volatile memory device with an oxygen-decorated graphene electrode

\jour Adv. Mater.

\yr 2011

\vol 23

\issue 16

\pages 1889--1893

\crossref{10.1002/adma.201004444}





\Bibitem{kol:bib10}

\by Zhu~Y., Li~L., Zhang~C., Casillas~G. et. al.

\paper A seamless three-dimensional carbon nanotube graphene hybrid material

\jour Nat. Commun.

\yr 2012

\vol 3

\papernumber  1225

\crossref{10.1038/ncomms2234}





\Bibitem{kol:bib11}

\by Kakade~B.~A., Pillai~V.~K., Late~D.~J., Chavan~P.~G. et. al.

\paper High current density, low threshold field emission from functionalized carbon nanotube bucky paper

\jour Appl. Phys. Lett.

\yr 2010

\vol 97

\issue 7

\papernumber  073102

\crossref{10.1063/1.3479049}



\Bibitem{kol:bib12}

\by Jousseaume~V., Cuzzocrea~J., Bernier~N., Renard~V.~T.

\paper Few Graphene layer/Carbon-Nanotube composite Grown at CMOS-compatible Temperature

\jour Appl. Phys. Lett.

\yr 2011

\vol 98

\papernumber 123103

\crossref{10.1063/1.3569142}





\Bibitem{kol:bib13}

\by Brenner~D.~W.

\paper Empirical potential for hydrocarbons for use in simulating the chemical vapor deposition of diamond films

\jour Phys. Rev. B

\yr 1990

\vol 42

\pages 9458--9471

\crossref{10.1103/PhysRevB.42.9458}





\Bibitem{kol:bib14}

\by Glukhova~O.~E., Kolesnikova~A.~S., Kossovich~E.~L., Zhnichkov~R.~Y.

\paper Super strong nanoindentors for biomedical applications based on bamboo-like nanotubes

\jour Proc. SPIE 

\vol 8233

\procinfo  2~February 2012

\volinfo Reporters, Markers, Dyes, Nanoparticles, and Molecular Probes for Biomedical Applications IV 

\yr 2012

\papernumber 823311

\crossref{10.1117/12.907035}





\Bibitem{kol:bib15}

\by Kolesnikova~A.~S.

\paper Mechanical properties of sorbents depending on nanopore sizes

\jour Phys. Sol. State

\yr 2017

\vol 59

\issue 7

\pages 1336--1339

\crossref{10.1134/S1063783417070113}





\Bibitem{kol:bib16}

\by Glukhova~O.~E., Kolesnikova~A.~S.

\paper Mechanical and emission properties of thinnest stable bamboolike nanotubes

\jour J. Phys. Conf. Ser.

\yr 2012

\vol 393

\papernumber 012027

\crossref{10.1088/1742-6596/393/1/012027}





\Bibitem{kol:bib17}

\by Lucas~A.~A., Lambin~P.~H., Smalley~R.~E.

\paper On the energetics of tubular fullerenes

\jour J. Phys. Chem. Solids

\yr 1993

\vol 54

\issue 5

\pages 587--593

\crossref{10.1016/0022-3697(93)90237-L}



\Bibitem{kol:bib18}

\by Glukhova~O.~E., Kolesnikova~A.~S., Slepchenkov~M.~M., Shmygin~D.~S.

\paper Atomic structure of energetically stable carbon nanotubes/graphene composites

\jour Phys. Sol. State

\yr 2015

\vol 57

\issue 5

\pages 1009--1013

\crossref{10.1134/S106378341505011X}



\Bibitem{kol:bib19}

\by Glukhova~O.~E., Slepchenkov~M.~M., Shmygin~D.~S.

\paper Nanoindentation of a new graphene/phospholipid composite: a numerical simulation

\jour Proc. SPIE 

\vol  10079

\volinfo Reporters, Markers, Dyes, Nanoparticles, and Molecular Probes for Biomedical Applications IX

\yr 2017

\procinfo 21~February 2017

\papernumber 1007910

\crossref{10.1117/12.2256809}





\RBibitem{kol:bib20}

\by Глухова~О.~Е., Гришина~О.~А., Савостьянов~Г.~В.

\paper Наноиндентирование липопротеинов высокой плотности углеродными нанотрубками: мультимасштабное моделирование

\jour Российский журнал биомеханики

\yr 2014

\vol 18

\issue 3

\pages 367--380

\elib{https://elibrary.ru/item.asp?id=22658811}

%\misctransl

%\lang In Russian

%\by Glukhova~O.~E., Grishina~O.~A., Savost'yanov~G.~V.

%\paper Nanoindentirovanie lipoproteinov vysokoj plotnosti uglerodnymi nanotrubkami: mul'timasshtabnoe modelirovanie 

%\jour Rossijskij zhurnal biomekhaniki

%\yr 2014

%\vol 18

%\issue 3

%\pages 367--380



\RBibitem{kol:bib21}

\by Глухова~О.~Е., Шунаев~В.~В.

\paper Исследование прочности на разрыв моно- и бислойного графена

\jour Нано- и микросистемная техника

\yr 2012

\issue 7(144)

\pages 25--29

\elib{https://elibrary.ru/item.asp?id=17847247}

%\misctransl

%\lang In Russian

%\by Glukhova~O.~E., Shunaev~V.~V.

%\paper Issledovanie prochnosti na razryv mono- i bislojnogo grafena

%\jour Nano- i mikrosistemnaya tekhnika

%\yr 2012

%\issue 7

%\pages 25--29



\end{thebibliography}





\newpage



\makeentitle





\footnotetext{\fullauthor{\enpid{73316}{Gaik~A. Baregamyan}}

%\orcid{0000-0000-0000-xxxx} % ORCID ID

\position{Postgraduate Student} 

\dept{Dept. of Mathematical Theory of Elasticity and Biomechanics}

\email{baregaik@gmail.com}

\fullauthor{\enpid{14470}{Leonid Yu. Kossovich}}

\orcid{0000-0002-4775-7348} % ORCID ID avtora 

\regalia{Dr. Phys. \& Math. Sci., Professor} 

\position{Head of Department} 

\dept{Dept. of Mathematical Theory of Elasticity and Biomechanics}

\email{president@sgu.ru}}



\vspace{-5mm}



\AdditionalContent{Competing interests}{We declare that we have no conflicts of interest in the authorship and publication of this article.} 

\AdditionalContent{Authors' contributions and responsibilities}{Each

author has participated in the article concept development and in

the manuscript writing. The authors are absolutely responsible for

submitting the final manuscript in print. Each author has approved

the final version of manuscript.}



\AdditionalContent{Funding}{This work was supported by Presidential scholarship SP-2502.2016.1.}







\begin{thebibliography}{10}





\Bibitem{kol:bib1:en}

\by Wei~J., Wei~C., Su~L., Fu~J., Lv~J. 

\yr 2015

\paper Synergistic Reinforcement of Phenol-Formaldehyde Resin Composites by Poly(Hexanedithiol)/Graphene Oxide

\jour  J.~Mater. Sci. Chem. Eng.

\vol 3

\issue 8

\pages 56--70

\crossref{10.4236/msce.2015.38009}



\Bibitem{kol:bib2:en}

\by Potts~J.~R., Dreyer~D.~R., Bielawski~C.~W., Ruoff~R.~S.

\paper Graphene-based polymer nanocomposites

\jour Polymer

\yr 2011

\vol 52

\issue 1

\pages 5--25

\crossref{10.1016/j.polymer.2010.11.042}





\Bibitem{kol:bib3:en}

\by Jung~N., Kwon~S., Lee~D., Yoon~D.~M. et. al.

\paper Synthesis of Chemically Bonded Graphene/ Carbon Nanotube Composites and their Application in Large Volumetric Capacitance Supercapacitors

\jour Adv. Mater.

\yr 2013

\vol 25

\issue 47

\pages 6854--6858

\crossref{10.1002/adma.201302788}



\Bibitem{kol:bib4:en}

\by Kim~S.~H., Song~W., Jung~M.~W., Kang~M.~A. et. al.

\paper Carbon nanotube and graphene hybrid thin film for transparent electrodes and field effect transistors

\jour Adv. Mater.

\yr 2014

\vol 26

\issue 25

\pages 4247--4252

\crossref{10.1002/adma.201400463}





\Bibitem{kol:bib5:en}

\by Kholmanov~I.~N., Magnuson~C.~W. , Piner~R., Kim~J.~Y. et. al.

\paper Optical, electrical, and electromechanical properties of hybrid graphene/carbon nanotube films

\jour Adv. Mater.

\yr 2015

\vol 27

\issue 19

\pages 3053--3059

\crossref{10.1002/adma.201500785}





\Bibitem{kol:bib6:en}

\by Dong~X., Li~B.,  Wei~A., Cao~X. et. al.

\paper One-step growth of graphene–carbon nanotube hybrid materials by chemical vapor deposition

\jour Carbon

\yr 2011

\vol 49

\issue 9

\pages 2944--2949

\crossref{10.1016/j.carbon.2011.03.009}



\Bibitem{kol:bib7:en}

\by Trist\'an-L\'opez~F., Morelos-G\'omez~A., Vega-Diaz~S.~M., Garcia-Betancourt~M.~L. et. al.

\paper Large area films of alternating graphene-carbon nanotube layers processed in water

\jour ACS Nano

\yr 2013

\vol 7

\issue 12

\pages 10788--10798

\crossref{10.1021/nn404022m}





\Bibitem{kol:bib8:en}

\by Mitrofanov~V.~V., Slepchenkov~M.~M., Zhang~G., Glukhova~O.~E.

\paper Hybrid Carbon Nanotube–Graphene Monolayer Films: Regularities of Structure, Electronic and Optical Properties

\jour Carbon

\yr 2017

\vol 115

\pages 803--810

\crossref{10.1016/j.carbon.2017.01.040}





\Bibitem{kol:bib9:em}

\by Yu~W.~J., Chae~S.~H., Lee~S.~Y., Duong~D.~L., Lee~Y.~H.

\paper Ultra-transparent, flexible single-walled carbon nanotube non-volatile memory device with an oxygen-decorated graphene electrode

\jour Adv. Mater.

\yr 2011

\vol 23

\issue 16

\pages 1889--1893

\crossref{10.1002/adma.201004444}





\Bibitem{kol:bib10:em}

\by Zhu~Y., Li~L., Zhang~C., Casillas~G. et. al.

\paper A seamless three-dimensional carbon nanotube graphene hybrid material

\jour Nat. Commun.

\yr 2012

\vol 3

\papernumber  1225

\crossref{10.1038/ncomms2234}





\Bibitem{kol:bib11:e,}

\by Kakade~B.~A., Pillai~V.~K., Late~D.~J., Chavan~P.~G. et. al.

\paper High current density, low threshold field emission from functionalized carbon nanotube bucky paper

\jour Appl. Phys. Lett.

\yr 2010

\vol 97

\issue 7

\papernumber  073102

\crossref{10.1063/1.3479049}



\Bibitem{kol:bib12?}

\by Jousseaume~V., Cuzzocrea~J., Bernier~N., Renard~V.~T.

\paper Few Graphene layer/Carbon-Nanotube composite Grown at CMOS-compatible Temperature

\jour Appl. Phys. Lett.

\yr 2011

\vol 98

\papernumber 123103

\crossref{10.1063/1.3569142}





\Bibitem{kol:bib13:E}

\by Brenner~D.~W.

\paper Empirical potential for hydrocarbons for use in simulating the chemical vapor deposition of diamond films

\jour Phys. Rev. B

\yr 1990

\vol 42

\pages 9458--9471

\crossref{10.1103/PhysRevB.42.9458}





\Bibitem{kol:bib14?e}

\by Glukhova~O.~E., Kolesnikova~A.~S., Kossovich~E.~L., Zhnichkov~R.~Y.

\paper Super strong nanoindentors for biomedical applications based on bamboo-like nanotubes

\jour Proc. SPIE 

\vol 8233

\procinfo  2~February 2012

\volinfo Reporters, Markers, Dyes, Nanoparticles, and Molecular Probes for Biomedical Applications IV 

\yr 2012

\papernumber 823311

\crossref{10.1117/12.907035}





\Bibitem{kol:bib15?E}

\by Kolesnikova~A.~S.

\paper Mechanical properties of sorbents depending on nanopore sizes

\jour Phys. Sol. State

\yr 2017

\vol 59

\issue 7

\pages 1336--1339

\crossref{10.1134/S1063783417070113}





\Bibitem{kol:bib16?E}

\by Glukhova~O.~E., Kolesnikova~A.~S.

\paper Mechanical and emission properties of thinnest stable bamboolike nanotubes

\jour J. Phys. Conf. Ser.

\yr 2012

\vol 393

\papernumber 012027

\crossref{10.1088/1742-6596/393/1/012027}





\Bibitem{kol:bib17?E}

\by Lucas~A.~A., Lambin~P.~H., Smalley~R.~E.

\paper On the energetics of tubular fullerenes

\jour J. Phys. Chem. Solids

\yr 1993

\vol 54

\issue 5

\pages 587--593

\crossref{10.1016/0022-3697(93)90237-L}



\Bibitem{kol:bib18?E}

\by Glukhova~O.~E., Kolesnikova~A.~S., Slepchenkov~M.~M., Shmygin~D.~S.

\paper Atomic structure of energetically stable carbon nanotubes/graphene composites

\jour Phys. Sol. State

\yr 2015

\vol 57

\issue 5

\pages 1009--1013

\crossref{10.1134/S106378341505011X}



\Bibitem{kol:bib19E?}

\by Glukhova~O.~E., Slepchenkov~M.~M., Shmygin~D.~S.

\paper Nanoindentation of a new graphene/phospholipid composite: a numerical simulation

\jour Proc. SPIE 

\vol  10079

\volinfo Reporters, Markers, Dyes, Nanoparticles, and Molecular Probes for Biomedical Applications IX

\yr 2017

\procinfo 21~February 2017

\papernumber 1007910

\crossref{10.1117/12.2256809}





\Bibitem{kol:bib20?E}

\jour Russian Journal of Biomechanics

\yr 2014

\vol 18

\issue 3

\pages 320--331

\paper Nanoindentation of high density lipoproteins by carbon nanotubes: a multiscale modelling

\by  Glukhova~O.~E. Grishina~O.~A., Savostyanov~G.~V.



\Bibitem{kol:bib21E?}

\lang In Russian

\by  Glukhova~O.~E., Shunaev~V.~V.

\paper Investigation of the tensile strength of mono-and bilayer graphene

\jour J.~Nano and Microsystem Technique

\yr 2012

\issue 7(144)

\pages 25--29



\end{thebibliography}



\renewcommand{\baselinestretch}{0.9}



%\end{document}

